% Generated draft: requires mathematical and source review.
\documentclass[11pt]{article}
\usepackage[margin=25mm]{geometry}
\usepackage{fontspec}
\setmainfont{Latin Modern Roman}
\usepackage{amsmath,amssymb,mathtools,unicode-math}
\setmathfont{Latin Modern Math}
\usepackage{enumitem,needspace,xurl,hyperref}
\setlength{\parindent}{0pt}
\setlength{\parskip}{5pt}
\setlength{\emergencystretch}{3em}
\newcommand{\N}{\mathbb{N}}
\newcommand{\Z}{\mathbb{Z}}
\newcommand{\Q}{\mathbb{Q}}
\newcommand{\R}{\mathbb{R}}
\newcommand{\C}{\mathbb{C}}
\newcommand{\F}{\mathbb{F}}
\newcommand{\A}{\mathbb{A}}
\newcommand{\PP}{\mathbb{P}}
\newcommand{\E}{\mathbb{E}}
\newcommand{\eps}{\varepsilon}
\newcommand{\dd}{\,\mathrm{d}}
\begin{document}
\section*{UnsolvedMath ID: 20003339}
\textbf{Source problem code:} catalogue entry
\subsection{Definitions and mathematical statement}
This record is the middle of the paragraph following Question 45, not an independent question. It predicts a normalized two-gap count tending to \((b-a)(d-c)\) for a ``generic'' quadratic form, but omits the definition of \(N_T\), the form's dimension, the full cutoff, and the meaning of generic. Its final sentence is unfinished. The aid reports that the next record completes a historical comment that density was then unknown, with the source typo ``Conjecture 45'' despite the heading ``Question 45.'' No complete arithmetic conjecture is reconstructed from this fragment.

The supplied normalization uses the pushforward density of Lebesgue measure on a cutoff \(D\) under \(Q\):
\[
 Q_*(1_D\,\dd x)=r_{Q,D}(t)\,\dd t,
 \qquad\kappa(Q,D)=\int_\R r_{Q,D}(t)^3\,\dd t.
\]
At regular values, the standard coarea interpretation is
\[
 r_{Q,D}(t)=\int_{D\cap Q^{-1}(t)}\frac{\dd\mathcal H^{d-1}(x)}{\|\nabla Q(x)\|}.
\]
The aid reports that the intended continuum and inhomogeneous-Poisson triple model forces normalization \(\kappa(Q,D)T^2\). The exact deterministic lattice-count definition and scaling conventions are absent, so this is retained as the model normalization, not a proof for the original quadratic-form count. It also reports an angular formula for the Euclidean unit-ball cutoff, without including that formula or its dimension/singularity conventions in this record.

A precise version of the supplied Poisson model is obtained from a nonzero bounded compactly supported density \(r\geq0\), with \(0<\kappa=\int r^3<\infty\), by giving a real-line Poisson process intensity \(r_T(s)=r(s/T^2)\). For bounded intervals \(I=(a,b)\), \(J=(c,d)\), let \(Z_T(I,J)\) count ordered triples of pairwise distinct process points \((s_0,s_1,s_2)\) with \(s_1-s_0\in I\), \(s_2-s_0\in J\). The reported continuum mean and overlap-variance conclusions are
\[
 \E Z_T\sim\kappa T^2|I||J|,\qquad
 \operatorname{Var}(Z_T)=O(T^2),
\]
\[
 \frac{Z_T}{\kappa T^2}\longrightarrow|I||J|\quad\text{in }L^2.
\]
Bounded compact support and this scaling are explicit hypotheses; no rate or conclusion for all unbounded or singular coarea densities is inferred. Nor is a Poisson assumption imposed on actual quadratic-form values.

The aid's one-point obstruction is the randomly shifted integer lattice \(\Z+U\), with \(U\) uniform on \([0,1)\). Its expected counting measure is exactly Lebesgue measure, but every difference of its points is an integer. Its two-gap count for
\[
 I=J=(1/4,3/4)
\]
is therefore zero, despite \(|I||J|=1/4\). This separates one-point intensity from the proposed two-gap law. It is a stochastic point-process countermodel, not an actual generic-quadratic-form counterexample. The normalization and model calculation remain a partial logical analysis of the unreproduced Question 45, not a resolution of its arithmetic distribution or density problem.
\subsection{Short English statement}
The fragment's quadratic-form count is not fully defined. The supplied continuum/Poisson model gives a cubic-density normalization and an \(L^2\) law under bounded-density hypotheses. A shifted integer lattice has the same one-point intensity but no indicated two-gap triples, so first-order information alone is insufficient.
\subsection{Sources}
\begin{itemize}
\item Local UnsolvedMath registry, permanent ID \(20003339\), exact source code \texttt{catalogue entry}; canonical \texttt{aim-other-notes.json}, \texttt{notes[145]}, AIM domain Other, registry category Miscellaneous, source tag \texttt{question}.
\item AIM workshop ``Emerging applications of measure rigidity,'' continuation after Question 45: \url{https://aimath.org/WWN/measrigid/measrigid.pdf}. The supplied verification says the words after ``set in'' lie in \texttt{AIM-OTHER-0147}, and ``Conjecture 45'' is a source-level cross-reference typo. The running page heading is not mathematical text. No PDF was fetched.
\item The original literature field is empty. No authoritative bibliographic citations, complete deterministic count, or explicit angular-density formula are supplied in this record.
\item Research attempt 1: result \texttt{context\_only}, problem status at run \texttt{not\_a\_problem}, rigor label \texttt{proved}, novelty confidence \texttt{low}; classification \texttt{OPEN-TRIAGE}, registry status \texttt{open}, suggested difficulty L2. Created and updated 14 August 2026. These are supplied assessments, not an arithmetic equidistribution theorem.
\item Attribution: Ulam AI / UnsolvedMath Contributors, CC BY 4.0. The machine-generated research aid is not peer reviewed and requires independent verification before citation.
\end{itemize}
\end{document}
